% Proposed replacement for the proof of Lemma \ref{thm:sKt_witnesses}.
% Mathematical proof verified via EPPAIII/CliqueCovers/AllParametersClassification.lean.
% Source manuscript intentionally unchanged.
\begin{proof}
Put $B=V(H)\setminus V(G)$. Since $H$ is regular and
$|B|=|G|$, counting edges between the two halves gives
\[
 |E(H[B])|=|E(G)|=s\binom t2.                    \tag{*}
\]
We first dispose of two simple cases. If $s=1$, then
$H[B]\cong K_t$. All vertices of $B$ have the same number
$k$ of neighbours in $G$, and the extensions of
$\Aut(G)=\Sym(t)$ show that $\binom tk\le t$.
Thus $k\in\{0,1,t-1,t\}$. In the middle two cases the
$t$ neighbourhoods are all distinct, giving a matching
or its complement; in the other cases the cross-edges
are empty or complete. If $t=1$, then $H[B]$ is empty
by~(*). Applying the same argument to the symmetric
action on the $s$ independent vertices of $G$ gives
the four cross-relations in the statement.

Now assume $s,t\ge2$. If $H$ is a disjoint union of
cliques, an edge between $C_i$ and $D_j$ forces
$C_i\cup D_j$ to be a clique. Since the two halves
are each partitioned into maximal cliques of size $t$,
the nonempty cross-relations form a matching between
the $C_i$ and $D_j$, and each is complete. By EPPA,
all permutations of the $C_i$ extend to automorphisms
of $H$, so this matching is either empty or perfect.
Both alternatives have the required connection matrix.

We may therefore suppose that $H$ is not a disjoint
union of cliques. For $w\in B$, record the profile
\[
 P(w)=(N_H(w)\cap C_1,\ldots,N_H(w)\cap C_s).
\]
Every permutation in
$\Aut(G)=\Sym(t)\wr\Sym(s)$ extends to an automorphism
of $H$ preserving $G$ and $B$. Thus the orbit of
$P(w)$ is realised by profiles of vertices of $B$,
and has size at most $st$. No coherent choice of
extensions is needed.

Call a coordinate \emph{partial} if its size is strictly
between $0$ and $t$, and let $p$ be the number of such
coordinates. If $2\le p<s$, the profile orbit has size at least
$\binom sp t^p>st$. If $p=s\ge2$, its size is at least
$t^s>st$, except when $s=t=2$. Hence there is at most
one partial coordinate, unless $s=t=2$.

Suppose that a profile has exactly one partial
coordinate of size $k$, and that $r$ of its other
coordinates are full. Its orbit has size
\[
  s\binom tk\binom{s-1}{r}\le st.
\]
It follows that $k\in\{1,t-1\}$ and
$r\in\{0,s-1\}$. In particular the orbit consists of
exactly $st$ distinct profiles and exhausts $B$.
The induced action on $B$ is then the natural action
of $\Sym(t)\wr\Sym(s)$ on $s$ blocks of size $t$.
Its orbits on unordered pairs are the pairs within
one block and those between different blocks.
Together with~(*), this forces $H[B]\cong sK_t$.
The partial coordinate gives a matching or its
complement on the diagonal, and all off-diagonal
entries are uniformly empty or complete.

If $s=t=p=2$, the four outside profiles are the
four choices of one vertex from each $C_i$.
Their pair orbits have sizes $4$ and $2$.
By~(*) the outside graph has two edges, so these
are the two opposite-corner edges. The resulting
graph is the Wagner graph, with the displayed matrix.

It remains to assume that no profile has a partial
coordinate. Then all vertices of each $C_i$ are
true twins in $H$. By transitivity their twin
classes have a common size $a\ge t$.
The vertex-transitive quotient on the twin classes
has at most $2s$ vertices and contains the $s$
selected classes as an independent set.
It has positive degree, since $H$ is not a disjoint
union of cliques. A regular graph of positive degree
has no independent set larger than half its order.
Consequently there are exactly $2s$ twin classes,
all of size $t$, and $B$ consists of the remaining
$s$ classes. Equation~(*) shows that no edges join
distinct classes contained in $B$.

Each outside class is connected completely to a
subset of the selected classes. These subsets have
a common size $r$ by regularity of $H$, and
permutations of the selected $C_i$ realise all
$\binom sr$ such subsets. Since there are only
$s$ outside classes, $\binom sr\le s$, whence
$r\in\{0,1,s-1,s\}$. For $r=1$ or $s-1$, pairing
outside classes with their distinguished $C_i$
gives a uniform diagonal and off-diagonal relation.
For $r=0$ or $s$, all cross-relations are equal.
This gives case~(1) and completes the proof.
\end{proof}
